\documentclass[11pt]{article}
\usepackage[margin=0.95in]{geometry}
\usepackage[T1]{fontenc}
\usepackage{lmodern,amsmath,amssymb,amsthm,mathtools,mathrsfs}
\usepackage{microtype,booktabs,longtable,array,enumitem,xcolor,fancyhdr,xurl}
\usepackage[colorlinks=true,linkcolor=blue!45!black,citecolor=blue!45!black,urlcolor=blue!50!black]{hyperref}
\usepackage{bookmark}
\setlength{\headheight}{14pt}
\pagestyle{fancy}\fancyhf{}
\fancyhead[L]{\small Triple Stieltjes correlations}
\fancyhead[R]{\small ProveIt research continuation}
\fancyfoot[C]{\thepage}
\emergencystretch=2em
\setcounter{tocdepth}{2}
\numberwithin{equation}{section}
\newtheorem{theorem}{Theorem}[section]
\newtheorem{proposition}[theorem]{Proposition}
\newtheorem{lemma}[theorem]{Lemma}
\newtheorem{corollary}[theorem]{Corollary}
\theoremstyle{definition}\newtheorem{definition}[theorem]{Definition}
\theoremstyle{remark}\newtheorem{remark}[theorem]{Remark}
\newtheorem{example}[theorem]{Example}
\DeclareMathOperator{\Li}{Li}
\DeclareMathOperator{\FP}{FP}
\DeclareMathOperator{\PV}{PV}
\DeclareMathOperator{\Res}{Res}
\DeclareMathOperator{\sgn}{sgn}
\DeclareMathOperator{\Log}{Log}
\newcommand{\ii}{\mathrm i}
\newcommand{\dd}{\,\mathrm d}
\newcommand{\TT}{\mathbb T}
\newcommand{\ZZ}{\mathbb Z}
\newcommand{\CC}{\mathbb C}
\newcommand{\RR}{\mathbb R}
\newcommand{\ee}[1]{\mathrm e^{2\pi\ii #1}}
\newcommand{\calB}{\mathcal B}
\newcommand{\calJ}{\mathcal J}
\newcommand{\calT}{\mathcal T}
\newcommand{\calM}{\mathcal M}
\newcommand{\Zjet}[2]{\zeta^{[#1]}(#2)}
\newcommand{\src}[1]{\nolinkurl{#1}}
\newcommand{\pin}{\texttt{8da6871fc2dd5762467320918ba23edd635bac4f}}
\title{\textbf{Triple Stieltjes Correlations}\\[5pt]
\large Colored Tornheim Jets, Harmonic Contact Terms,\\
and Exact Gamma Reductions}
\author{Research continuation prepared for Vladimir Reshetnikov\\
\small ProveIt programme \quad\textbullet\quad Prepared with OpenAI ChatGPT}
\date{10 October 2026}
\hypersetup{pdftitle={Triple Stieltjes Correlations: Colored Tornheim Jets, Harmonic Contact Terms, and Exact Gamma Reductions},pdfauthor={Prepared for Vladimir Reshetnikov with OpenAI ChatGPT}}
\begin{document}
\maketitle
\begin{abstract}
The shifted Hurwitz-jet report in the ProveIt repository asks whether the
canonical product of three separated Stieltjes singularities admits an
all-index description by colored Tornheim or multiple Lerch jets. We give
an affirmative answer: six colored Tornheim functions, each entire in its
three spectral variables when its two colors are nontrivial, generate every
such finite part. A second formula replaces their jets by three absolutely
convergent logarithmic Mellin integrals of polylogarithm order derivatives,
with an explicit endpoint constant at every index. This includes a fully
specified formula for three digamma factors at arbitrary distinct real,
and hence rational, shifts.

The same kernel gives all three-factor mean-zero primitive products. Complete homogeneous
harmonic numbers describe the primitives, while elementary harmonic numbers
describe the precise contact correction between periodic differentiation
and pointwise finite-part extension. A reflected subclass admits a stronger
reduction: one Stieltjes function multiplied by arbitrarily many separated
cotangent kernels reduces to finitely many single Stieltjes values. At
index zero and rational shifts, these become finite log-Gamma evaluations.
An explicit cubic principal value is
$\gamma+4\log2+3\log\pi-4\log\Gamma(1/4)$.

The proofs are analytic and algebraic. The accompanying implementation passes
621 exact finite assertions and 56 numerical comparisons, including three
independent evaluations of a genuinely singular three-digamma product.
Floating-point checks are diagnostics, not interval certificates. Classical
Fourier, partial-fraction, and Gamma identities are credited; global priority
for the derived families is not asserted.
\end{abstract}
\vspace{5pt}
\noindent\textbf{Claim boundary.} This article resolves the specified
three-shift representation question. It does not prove the remaining
$S_6$ or $S_8$ special-value conjectures, period independence, or minimal
transcendental depth. No repository files were modified and no
proof-assistant formalization is claimed.
\clearpage
\tableofcontents
\clearpage
% BEGIN sections/01-scope.tex
\section{Research question, scope, and antecedents}\label{tsc:scope}
The repository's continuation
\src{stieltjes-correlation-calculus/continuations/shifted-hurwitz-jets/}
ends with a concrete question about three distinct translations
\cite{tsc:shj}. Its bilinear theory has only one surviving Fourier index;
three factors impose one relation among three nonzero frequencies and leave
a genuine double sum. The proposed target is an all-index coordinate class,
not an unproved assertion that every resulting period reduces to one-variable
Stieltjes constants.

We resolve that target with two interchangeable representations. The first
is a finite coefficient formula in six colored Tornheim functions
(Theorems~\ref{tsc:master} and \ref{tsc:stieltjes-triple}). The second is an
explicit convergent integral formula using only spectral derivatives of
classical polylogarithms and a finite endpoint correction
(Theorem~\ref{tsc:mellin-all}). Neither definition invokes an unspecified
regularization, an integer-relation search, or a numerical limiting value.
The entire continuation of each colored double sum is proved before any
spectral coefficient is extracted.

The reflected subclass behaves differently from the generic triple.
Cotangent partial fractions force a finite reduction at every number of
separated factors. Combining that algebra with a normalized Stieltjes
transform gives Theorem~\ref{tsc:all-depth}, and rational shifts give
Theorem~\ref{tsc:rational}. The distinction is substantive: the generic
three-factor theorem supplies a colored depth-two representation; the
reflected theorem supplies a finite single-function reduction. The latter
does not imply the former has such a reduction.

\subsection{What is inherited and what is developed here}
The Hurwitz Fourier expansion, the shift relation, Lerch's formula at zero,
and the polygamma reflection law are classical \cite{tsc:dlmf-zeta,tsc:dlmf-gamma}.
Espinosa and Moll develop Hurwitz integral and primitive calculus
\cite{tsc:em1,tsc:em2}. Zhao gives finite double-polylogarithm reductions of
integer colored Tornheim sums \cite{tsc:zhao}. The unshifted cubic
log-Gamma integral already has a Tornheim-derivative evaluation due to
Bailey, D.~Borwein, and J.~M.~Borwein \cite{tsc:bbb}; the canonical ProveIt
chapter explicitly credits it \cite{tsc:gamma-chapter}. We do not claim any
of these results as discoveries of this article.

The repository already has an entire periodic Hurwitz family, bilinear
Stieltjes closure, contact corrections, and normalized primitive identities
\cite{tsc:repo,tsc:shj}. We reprove the normalization needed below so that
the trilinear statement is self-contained. The principal addition is the
separated three-factor proof chain: entire colored continuation, the six
frequency sectors, all-index extraction, convergent Mellin-jet coordinates,
and the resulting primitive and contact laws. The all-depth reflected
formula and its explicit evaluations are accompanying consequences, with
no claim of global priority for their individual specializations.

\subsection{Source snapshot and limits of the audit}
The final observed repository head was
\begin{center}\pin.\end{center}
The key shifted-jet source was read at that revision and has blob identifier
\src{990369336bdf2794aae3057c3231eed4ecadbef3}.
Other selected manuscript reads and the incoming-archive inventory are
recorded in \src{notes/source_snapshot.json}; not every earlier read was
pinned to the final head. A Library search also recovered the same
three-shift question and the abstract of the later nested-harmonic report.
The seven incoming ZIP files were inventoried, not unpacked or analytically
audited in full. This is a targeted continuation, not an audit of every
archived claim.

The current manuscript records a proof of $S_4$, while $S_6$ and the revised
$S_8$ candidate remain conjectural \cite{tsc:repo}. The present argument does
not change those statuses. It also does not turn a representation theorem
into a proof that its coordinates are arithmetically independent.

% END sections/01-scope

% BEGIN sections/02-normalization.tex
\section{Canonical endpoint normalization}\label{tsc:normalization}
Write $\TT=\RR/\ZZ$ and $e(x)=\exp(2\pi\ii x)$. Fourier coefficients use
\[
 \widehat f(n)=\langle f,e(-nx)\rangle.
\]
For $n\ne0$ fix
\begin{equation}\label{tsc:branch}
 (2\pi\ii n)^{-u}=(2\pi|n|)^{-u}
       \exp\!\left(-\frac{\pi\ii u}{2}\sgn n\right).
\end{equation}
All shifts below are real modulo one. Values assigned at isolated seam
points do not affect ordinary integrals. Spectral derivatives are denoted
$\zeta^{[j]}(s,x)=\partial_s^j\zeta(s,x)$ and
$\Li_s^{[j]}(z)=\partial_s^j\Li_s(z)$. Parenthesized superscripts on
$\gamma_m$ or $\psi$ denote derivatives in the spatial argument.

We deliberately use the spectral coordinate $u=1-s$:
\begin{equation}\label{tsc:stieltjes-convention}
 \zeta(1-u,x)=-\frac1u+
     \sum_{m\ge0}\frac{\gamma_m(x)}{m!}u^m,
 \qquad \gamma_0(x)=-\psi(x).
\end{equation}
This changes no Stieltjes convention: it is the usual Laurent expansion
at $s=1$ after substituting $s=1-u$. The shift relation gives
\begin{equation}\label{tsc:gamma-local}
 \gamma_m(x)=\frac{\log^m x}{x}+\gamma_m(1+x),\qquad x>0.
\end{equation}

\begin{definition}[Scale-one periodic Stieltjes extension]\label{tsc:Cdef}
For a smooth periodic test function $\varphi$, define
\begin{equation}\label{tsc:Cfp}
 \langle C_m,\varphi\rangle=
 \lim_{\varepsilon\downarrow0}\left\{
   \int_{\varepsilon}^{1}\gamma_m(x)\varphi(x)\dd x
    +\frac{\varphi(0)}{m+1}\log^{m+1}\varepsilon\right\}.
\end{equation}
The coordinate at the singularity is exactly $x$, without a rescaling or
an additional subtraction length. Set $C_{m,a}(x)=C_m(x-a)$.
\end{definition}
Existence follows from \eqref{tsc:gamma-local}: after subtracting
$\varphi(0)\log^m x/x$, the integrand is $O(1+|\log x|^m)$.

\begin{lemma}[Entire normalized Hurwitz family]\label{tsc:entire-B}
The Fourier prescription
\begin{equation}\label{tsc:Bdef}
 \calB_u(x)=\sum_{n\ne0}(2\pi\ii n)^{-u}e(nx)
\end{equation}
defines an entire distribution-valued function of $u\in\CC$, with mean
zero and singular support contained in $0$. On $0<x<1$,
\begin{equation}\label{tsc:Bhurwitz}
 \calB_u(x)=\frac{\zeta(1-u,x)}{\Gamma(u)},
\end{equation}
with removable values interpreted analytically. Near $u=0$,
\begin{equation}\label{tsc:Zlaurent}
 \Gamma(u)\calB_u=\frac{\delta_0-1}{u}
          +\sum_{m\ge0}\frac{C_m}{m!}u^m.
\end{equation}
In particular, $\calB_0=\delta_0-1$ and $\langle C_m,1\rangle=0$.
\end{lemma}
\begin{proof}
On a compact set of spectral parameters, the coefficients in
\eqref{tsc:Bdef}, and every spectral derivative of them, grow at most
polynomially in $|n|$, with additional powers of $\log|n|$. The Fourier
coefficients of a smooth test function decrease faster than every power.
Their pairing is therefore normally convergent, including all spectral
derivatives. This proves entire dependence in distributions.

For $\Re u>1$, the classical Hurwitz Fourier expansion
\cite[Eq.~25.11.9]{tsc:dlmf-zeta} gives \eqref{tsc:Bhurwitz}; both sides
continue analytically away from the seam. The right side is smooth there
and entire in $u$, since the pole of $\zeta(1-u,x)$ at zero is canceled by
$1/\Gamma(u)$. This also proves the assertion about singular support.

For the Laurent normalization, split a test-function integral near zero
and use
\[
 \zeta(1-u,x)=x^{u-1}+\zeta(1-u,1+x).
\]
The first term has residue $\varphi(0)$ at $u=0$; the second has residue
$-\int\varphi$. The regular coefficients of the first term are precisely
the scale-one finite parts of $\log^m x/x$. The regular coefficients of
the second are smooth. This proves \eqref{tsc:Zlaurent} and the stated
normalization. Mean zero follows either from the Fourier prescription or
by comparing Laurent coefficients of its zero Fourier coefficient.
\end{proof}

\subsection{Products at separated singularities}
Let $a_1,a_2,a_3$ be pairwise distinct in $\TT$. Near $a_i$, the other two
translated factors are smooth. A partition of unity therefore defines a
unique product of these distributions: multiply the singular factor by the
smooth factors on each such neighborhood and use ordinary multiplication
elsewhere. The definition is independent of the partition. It also gives
joint entire dependence for products of the $\calB_{u_i}$, because locally
only one distribution is multiplied by smooth parameter-dependent factors.
No assertion about products at a common singularity is needed.

\begin{definition}[Separated triple finite part]\label{tsc:Jdef}
For $\mathbf m=(m_1,m_2,m_3)\in\ZZ_{\ge0}^3$, put
\begin{equation}\label{tsc:J}
 \calJ_{\mathbf m}(\mathbf a)=
  \left\langle\prod_{i=1}^3 C_{m_i,a_i},1\right\rangle
  =\FP\int_0^1\prod_{i=1}^3
       \gamma_{m_i}(\{x-a_i\})\dd x.
\end{equation}
At each seam $a_i$, omit the right-hand arc $0<t<\varepsilon_i$ and add
\begin{equation}\label{tsc:triple-counterterm}
 \frac{\log^{m_i+1}\varepsilon_i}{m_i+1}
       \prod_{j\ne i}\gamma_{m_j}(\{a_i-a_j\}).
\end{equation}
The three cutoffs may tend to zero independently.
\end{definition}
The product of the two smooth factors differs from its seam value by
$O(t)$, so the same local subtraction proof establishes existence. The
one-sided nature of these counterterms must be distinguished from the
symmetric principal values used later for cotangent poles.

% END sections/02-normalization

% BEGIN sections/03-tornheim.tex
\section{Colored Tornheim functions and entire continuation}\label{tsc:tornheim}
For an initially convergent double sum, write
\begin{equation}\label{tsc:Tdef}
 T(r,s,t;z,w)=\sum_{m,n\ge1}
           \frac{z^m w^n}{m^r n^s(m+n)^t}.
\end{equation}
For example, $\Re r,\Re s,\Re t>1$ is more than sufficient for unit
colors. The functions below are not defined by assigning a value to a
divergent instance of this series. They are defined by its unique analytic
continuation, constructed explicitly in the next theorem.

\begin{theorem}[Nontrivial-color entireness]\label{tsc:T-entire}
Fix $|z|=|w|=1$ with $z\ne1$ and $w\ne1$. Then
$T(r,s,t;z,w)$ extends to a jointly entire function of $(r,s,t)\in\CC^3$.
For $\Re t>0$ and arbitrary $r,s\in\CC$ it is given by
\begin{equation}\label{tsc:T-mellin}
 T(r,s,t;z,w)=\frac1{\Gamma(t)}
   \int_0^\infty x^{t-1}\Li_r(ze^{-x})\Li_s(we^{-x})\dd x.
\end{equation}
The two colors may coincide. At every nonnegative integer $k$,
\begin{equation}\label{tsc:T-negative}
 T(r,s,-k;z,w)=\sum_{j=0}^k\binom{k}{j}
                \Li_{r-j}(z)\Li_{s-k+j}(w).
\end{equation}
\end{theorem}
\begin{proof}
Put $F(x)=\Li_r(ze^{-x})\Li_s(we^{-x})$. Because neither color is one,
the arguments have neighborhoods avoiding the polylogarithm singularity
at one. The standard analytic continuation of $\Li_r$ is entire in its
order there \cite{tsc:dlmf-poly}. Thus $F$ is analytic in $x$ near zero,
jointly with $r,s$. On compact sets of orders, the power series for large
$x$ gives an exponentially decreasing bound for $F$ and its spectral
derivatives. These facts prove convergence and holomorphy of
\eqref{tsc:T-mellin} for $\Re t>0$. Expanding the two polylogarithms where
absolute convergence permits and integrating $e^{-(m+n)x}$ identifies it
with \eqref{tsc:Tdef} on a nonempty open set.

The Taylor coefficients at the lower endpoint are
\begin{equation}\label{tsc:Fcoeff}
 f_k=[x^k]F(x)=\frac{(-1)^k}{k!}
        \sum_{j=0}^k\binom{k}{j}
                   \Li_{r-j}(z)\Li_{s-k+j}(w).
\end{equation}
For any positive integer $N$, the expression
\begin{align}\label{tsc:T-subtracted}
 \frac1{\Gamma(t)}\bigg\{&
  \int_0^1 x^{t-1}\left(F(x)-\sum_{k=0}^{N-1}f_kx^k\right)\dd x
  +\int_1^\infty x^{t-1}F(x)\dd x
  +\sum_{k=0}^{N-1}\frac{f_k}{t+k}\bigg\}
\end{align}
is holomorphic for $\Re t>-N$ after removing the apparent singularities
at $0,-1,\ldots,-N+1$. Local uniformity in $r,s,t$ follows from the Taylor
remainder and exponential decay. The only poles of the braces are the
listed simple ones, and the simple zeros of $1/\Gamma(t)$ cancel them.
The expressions for different $N$ agree on overlaps, which proves joint
entireness. At $t=-k$, the zero of $1/\Gamma(t)$ has leading coefficient
$(-1)^k k!$; multiplication by the residue $f_k$ gives
\eqref{tsc:T-negative}.
\end{proof}

\begin{remark}[Why the color assumption matters]
At $z=1$ or $w=1$, the Taylor argument at $x=0$ changes: a polylogarithm
can have a fractional-power or logarithmic local term. The theorem does
not assert global entireness in that case. Later primitive formulas can
allow color one near positive spectral integers because the original
Tornheim series then converges absolutely. These are different reasons
for analyticity, and they must not be conflated.
\end{remark}

\subsection{An explicit jet at the zero third order}
For fixed nontrivial colors, define ordinary convergent logarithmic moments
\begin{align}\label{tsc:Mmoments}
 M_j(r,s;z,w)={}&\int_0^1\frac{\log^j x}{x}
        \bigl(\Li_r(ze^{-x})\Li_s(we^{-x})-\Li_r(z)\Li_s(w)\bigr)\dd x\notag\\
 &+\int_1^\infty\frac{\log^j x}{x}
                    \Li_r(ze^{-x})\Li_s(we^{-x})\dd x.
\end{align}
Taking $N=1$ in \eqref{tsc:T-subtracted} yields, locally for $|t|<1$,
\begin{equation}\label{tsc:T-zero-jet}
 T(r,s,t;z,w)=\frac1{\Gamma(1+t)}
  \left\{\Li_r(z)\Li_s(w)+
               \sum_{j\ge0}\frac{M_j(r,s;z,w)}{j!}t^{j+1}\right\}.
\end{equation}
Every derivative in $r,s$ may be passed through the moments. In particular,
\[
 T(r,s,0;z,w)=\Li_r(z)\Li_s(w),\qquad
 T_t(r,s,0;z,w)=\gamma\Li_r(z)\Li_s(w)+M_0(r,s;z,w).
\]
These formulas specify the third-order jets by convergent integrals rather
than by a limiting procedure at a pole.

\subsection{Classical double-polylogarithm and harmonic coordinates}
Use the convention
\[
 \Li_{p,q}(z,w)=\sum_{N>n\ge1}\frac{z^Nw^n}{N^p n^q}
\]
where convergent, followed by analytic continuation when justified.
Changing variables $N=m+n$ gives
\begin{equation}\label{tsc:T-zero-first}
 T(0,s,t;z,w)=\Li_{t,s}(z,w/z).
\end{equation}
For $z\ne w$, summing the finite geometric progression at fixed $m+n$
also gives
\begin{equation}\label{tsc:T00}
 T(0,0,t;z,w)=\frac{w\Li_t(z)-z\Li_t(w)}{z-w}.
\end{equation}
For positive integers $r,s$, the standard partial-fraction reduction
underlying Zhao's result \cite{tsc:zhao} is
\begin{align}\label{tsc:integer-reduction}
 T(r,s,t;z,w)={}&\sum_{j=1}^r\binom{r+s-j-1}{s-1}
                     \Li_{t+r+s-j,j}(w,z/w)\notag\\
 &+\sum_{j=1}^s\binom{r+s-j-1}{r-1}
                     \Li_{t+r+s-j,j}(z,w/z).
\end{align}
For completeness, the rational identity being summed is
\[
 \frac1{m^r n^s}=
  \sum_{j=1}^r\frac{\binom{r+s-j-1}{s-1}}{m^j(m+n)^{r+s-j}}
 +\sum_{j=1}^s\frac{\binom{r+s-j-1}{r-1}}{n^j(m+n)^{r+s-j}}.
\]
It follows by repeated use of
$1/(mn)=1/(m(m+n))+1/(n(m+n))$, or by comparing principal parts in
$m$ for fixed $m+n$. Multiplying by $(m+n)^{-t}z^mw^n$ and summing proves
\eqref{tsc:integer-reduction} first in an absolutely convergent domain.

At unit arguments and in a convergence domain,
\[
 \Li_{p,q}(1,1)=\sum_{N\ge1}\frac{H_{N-1}^{(q)}}{N^p},
 \qquad H_N^{(q)}=\sum_{j=1}^N j^{-q}.
\]
Thus the colored coordinate class includes generalized harmonic Euler sums.
The integer reduction is a known identity, not a new theorem here.
Moreover, differentiating it in the freely complex variable $t$ is
legitimate; differentiating it in the discrete integers $r$ or $s$ is not
licensed by the formula. The whole point of Theorem~\ref{tsc:T-entire}
is to supply a genuine analytic family for those missing jets.

% END sections/03-tornheim

% BEGIN sections/04-triple.tex
\section{The six-sector kernel and the triple coefficient theorem}\label{tsc:triple}
Let $\mathbf s=(s_1,s_2,s_3)$ and let $\mathbf a=(a_1,a_2,a_3)$ have
pairwise distinct entries modulo one. For each $k\in\{1,2,3\}$, let $i,j$
be the other two indices in increasing order and put
\begin{equation}\label{tsc:colors}
 z_{ik}=e(a_k-a_i),\qquad z_{jk}=e(a_k-a_j).
\end{equation}
Define
\begin{align}\label{tsc:kernel}
 \calT(\mathbf s;\mathbf a)=\sum_{k=1}^3\bigg[&
  e^{-\pi\ii(s_i+s_j-s_k)/2}
       T(s_i,s_j,s_k;z_{ik},z_{jk})\notag\\
 &+e^{\pi\ii(s_i+s_j-s_k)/2}
       T(s_i,s_j,s_k;z_{ik}^{-1},z_{jk}^{-1})\bigg].
\end{align}
All six terms are entire in $\mathbf s$ by Theorem~\ref{tsc:T-entire}.
For real spectral variables the terms in each bracket are conjugate, but
\eqref{tsc:kernel} itself is the complex-analytic definition.

\begin{theorem}[Six-sector product identity]\label{tsc:master}
For all $\mathbf s\in\CC^3$,
\begin{equation}\label{tsc:master-eq}
 \left\langle\prod_{i=1}^3\calB_{s_i}(x-a_i),1\right\rangle
       =(2\pi)^{-s_1-s_2-s_3}\calT(\mathbf s;\mathbf a).
\end{equation}
The left side is the separated distribution product. For $\Re s_i>0$
it is an ordinary integral. In particular,
\begin{equation}\label{tsc:T000-check}
 \calT(0,0,0;\mathbf a)=2.
\end{equation}
\end{theorem}
\begin{proof}
First take $\Re s_i>1$, so all three Fourier series are absolutely
convergent. Integration leaves exactly the triples of nonzero integers
$(n_1,n_2,n_3)$ with $n_1+n_2+n_3=0$. Such a triple has either two
positive entries and one negative entry, or the reverse. If $k$ is the
negative singleton, write $n_i=m$, $n_j=n$, $n_k=-(m+n)$ with $m,n\ge1$.
The Fourier powers contribute
\[
 (2\pi)^{-s_1-s_2-s_3}
   e^{-\pi\ii(s_i+s_j-s_k)/2}
          m^{-s_i}n^{-s_j}(m+n)^{-s_k}.
\]
The translation factors contribute
$e(m(a_k-a_i)+n(a_k-a_j))=z_{ik}^m z_{jk}^n$.
Reversing all three signs produces the other term in the bracket. The
six disjoint sectors exhaust the Fourier constraint and give
\eqref{tsc:master-eq}.

Both sides are entire in all three variables. For the left side this was
proved locally at separated singularities; for the right side it follows
from the colored Mellin theorem. The identity theorem, applied successively
to the three variables, extends the equality to all of $\CC^3$.
For $\Re s_i>0$, the only singularity near $a_i$ is an integrable
power $t^{s_i-1}$, so the stated ordinary-integral interpretation holds.
Finally, separated delta functions have zero pairwise products, and
\[
 (\delta_{a_1}-1)(\delta_{a_2}-1)(\delta_{a_3}-1)
    =\delta_{a_1}+\delta_{a_2}+\delta_{a_3}-1.
\]
Pairing with one gives \eqref{tsc:T000-check}.
\end{proof}

\subsection{All Stieltjes indices}
Set
\begin{equation}\label{tsc:E}
 A=\gamma+\log(2\pi),\qquad
 E(u)=\Gamma(1+u)(2\pi)^{-u}
 =\exp\left(-Au+\sum_{k\ge2}\frac{(-1)^k\zeta(k)}k u^k\right).
\end{equation}
The power series is used near zero. Its coefficients are explicit Bell
polynomials in $A,\zeta(2),\ldots$.

\begin{theorem}[All-index triple Stieltjes closure]\label{tsc:stieltjes-triple}
For every nonnegative triple $\mathbf m$ and every separated real shift
triple $\mathbf a$,
\begin{equation}\label{tsc:all-index}
 \boxed{\displaystyle
 \calJ_{\mathbf m}(\mathbf a)=
  \left(\prod_{i=1}^3m_i!\right)
  [u_1^{m_1+1}u_2^{m_2+1}u_3^{m_3+1}]
       \left(\prod_{i=1}^3E(u_i)\right)\calT(\mathbf u;\mathbf a).}
\end{equation}
Every coefficient on the right is a finite combination of jets of the six
entire colored Tornheim functions. No diagonal Laurent finite part or
unspecified order of limits is present.
\end{theorem}
\begin{proof}
Multiply \eqref{tsc:master-eq} by $\prod_i\Gamma(u_i)$ and use
\eqref{tsc:Zlaurent}. The coefficient of
$u_1^{m_1}u_2^{m_2}u_3^{m_3}$ on the left is
$\calJ_{\mathbf m}/\prod_i m_i!$. Indeed, a residue factor in coordinate
$i$ has exponent $-1$ in that coordinate and no dependence on $u_i$ in
any other factor; it cannot contribute to a nonnegative exponent there.
On the right, $\Gamma(u_i)=\Gamma(1+u_i)/u_i$. Moving the three simple
coordinate denominators into the coefficient extraction gives
\eqref{tsc:all-index}. This also proves directly that simultaneous or
iterated coefficient extraction has the same meaning here.
\end{proof}

\subsection{The fully regularized three-digamma formula}
Define $\mathscr D_A=\prod_{i=1}^3(\partial_{s_i}-A)$. Since each
$\gamma_0=-\psi$, the first case of \eqref{tsc:all-index} is
\begin{equation}\label{tsc:digamma-operator}
 \FP\int_0^1\prod_{i=1}^3\psi(\{x-a_i\})\dd x
       =-\left.\mathscr D_A\calT(\mathbf s;\mathbf a)\right|_{\mathbf s=0}.
\end{equation}
This is already a finite, fully normalized answer to the first concrete
target in \cite{tsc:shj}. An equivalent three-term real form is useful.
Put $\alpha=A+\pi\ii/2$. Then
\begin{align}\label{tsc:digamma-three}
 \FP\int_0^1\prod_{i=1}^3\psi(\{x-a_i\})\dd x
 =-2\Re\sum_{k=1}^3
 \left.(\partial_r-\alpha)(\partial_s-\alpha)
       (\partial_t-\overline\alpha)
       T(r,s,t;z_{ik},z_{jk})\right|_{r=s=t=0}.
\end{align}
Only derivatives of order at most one in each of the three spectral
variables occur. The next section turns these derivatives into ordinary
convergent integrals.

\subsection{Translation and spectral lowering}
The Fourier definition gives $\partial_x\calB_s=\calB_{s-1}$ as an
entire distributional identity. Differentiating the separated product yields
\begin{equation}\label{tsc:lowering}
 \partial_{a_j}\calT(\mathbf s;\mathbf a)
       =-2\pi\calT(\mathbf s-\mathbf e_j;\mathbf a),
 \qquad
 \sum_{j=1}^3\calT(\mathbf s-\mathbf e_j;\mathbf a)=0.
\end{equation}
The second equality also follows from simultaneous translation invariance.
These are identities for the complete six-sector combination, not a license
to discard sectors or their phases.

% END sections/04-triple

% BEGIN sections/05-mellin-jets.tex
\section{Three convergent polylogarithmic integrals at every index}\label{tsc:polylog-moments}
The coefficient formula has a particularly explicit integral counterpart.
Let
\begin{equation}\label{tsc:cdalpha}
 \ell=\log(2\pi),\qquad c=\ell+\frac{\pi\ii}{2},\qquad
 d=\ell-\frac{\pi\ii}{2},\qquad \alpha=\gamma+c.
\end{equation}
For a nontrivial unit color $z$ and $x\ge0$, define
\begin{equation}\label{tsc:Vdef}
 V_m(z,x)=m![u^{m+1}]\,
      \Gamma(1+u)e^{-cu}\Li_u(ze^{-x}).
\end{equation}
This is a finite expression in ordinary polylogarithm order derivatives.
If
\[
 E_+(u)=\Gamma(1+u)e^{-cu}=\sum_{j\ge0}e_j^+u^j,
\]
then
\begin{equation}\label{tsc:Vfinite}
 V_m(z,x)=m!\sum_{r=0}^{m+1}\frac{e_{m+1-r}^+}{r!}
                         \Li_0^{[r]}(ze^{-x}),
\end{equation}
where
\begin{equation}\label{tsc:e-recurrence}
 e_0^+=1,\qquad
 n e_n^+=-\alpha e_{n-1}^+
          +\sum_{k=2}^n(-1)^k\zeta(k)e_{n-k}^+.
\end{equation}
For example, with $Z_r=\Li_0^{[r]}(ze^{-x})$,
\begin{align}\label{tsc:Vexamples}
 V_0&=Z_1-\alpha Z_0,\notag\\
 V_1&=\frac{Z_2}{2}-\alpha Z_1+
                  \frac{\alpha^2+\zeta(2)}2Z_0,\notag\\
 V_2&=\frac{Z_3}{3}-\alpha Z_2+
      (\alpha^2+\zeta(2))Z_1
      -\frac{\alpha^3+3\alpha\zeta(2)+2\zeta(3)}3 Z_0.
\end{align}
In particular, $Z_0(q)=q/(1-q)$ is elementary.

For $F$ analytic at zero and exponentially decreasing at infinity, put
\begin{align}\label{tsc:Mfunctional}
 \calM_{n,d}[F]={}&
    \int_0^1\frac{F(x)-F(0)}x(\log x-d)^n\dd x
   +\int_1^\infty\frac{F(x)}x(\log x-d)^n\dd x\notag\\
 &+\frac{(-d)^{n+1}}{n+1}F(0).
\end{align}
Both displayed integrals are ordinary absolutely convergent integrals.
The last term is essential; it is the contribution of the simple Mellin
pole after multiplication by $e^{-du}$.

\begin{theorem}[All-index convergent Mellin formula]\label{tsc:mellin-all}
With the colors \eqref{tsc:colors} and the conventions above,
\begin{equation}\label{tsc:M-all}
 \boxed{\displaystyle
 \calJ_{\mathbf m}(\mathbf a)=2\Re\sum_{k=1}^3
 \calM_{m_k,d}\!\big[
       V_{m_i}(z_{ik},\cdot)V_{m_j}(z_{jk},\cdot)\big].}
\end{equation}
Thus every separated triple Stieltjes finite part has a representation by
three convergent logarithmic integrals and three explicit endpoint terms.
Only finitely many polylogarithm order derivatives at order zero occur.
\end{theorem}
\begin{proof}
Consider the positive-frequency sector with singleton $k$. Before Laurent
coefficient extraction, multiply its Tornheim term by the three Gamma
factors appearing in \eqref{tsc:Zlaurent}. Inserting
\eqref{tsc:T-mellin} cancels the Gamma factor belonging to the singleton.
The sector becomes the analytic continuation of
\begin{equation}\label{tsc:sector-mellin}
 \Gamma(u_i)e^{-cu_i}\Gamma(u_j)e^{-cu_j}e^{-du_k}
 \int_0^\infty x^{u_k-1}
           \Li_{u_i}(z_{ik}e^{-x})\Li_{u_j}(z_{jk}e^{-x})\dd x.
\end{equation}
Taking the nonnegative Laurent coefficients in $u_i,u_j$, with the
factorials prescribed by \eqref{tsc:stieltjes-convention}, replaces the two
polylogarithmic factors by $V_{m_i}$ and $V_{m_j}$.

For the remaining variable $u=u_k$, write the Mellin integral as
\[
 \frac{F(0)}u+
 \int_0^1x^{u-1}(F(x)-F(0))\dd x+
 \int_1^\infty x^{u-1}F(x)\dd x.
\]
The coefficient $m_k![u^{m_k}]$ after multiplication by $e^{-du}$ is exactly
\eqref{tsc:Mfunctional}. The nontrivial colors ensure $F(x)-F(0)=O(x)$;
all of its spectral derivatives have the same property. At infinity they
decay exponentially. These local bounds justify every finite spectral
extraction under the subtracted integrals. Reversing the three Fourier
signs conjugates this expression because the shifts and Stieltjes indices
are real. Summing the three positive sectors and their conjugates proves
\eqref{tsc:M-all}.
\end{proof}

\begin{corollary}[A convergent three-digamma formula]\label{tsc:psi-mellin}
Put
\[
 W_z(x)=\left.\partial_s\Li_s(ze^{-x})\right|_{s=0}
            -\left(A+\frac{\pi\ii}{2}\right)\frac{ze^{-x}}{1-ze^{-x}}.
\]
Then
\begin{align}\label{tsc:psi-integrals}
 \FP\int_0^1\prod_{i=1}^3\psi(\{x-a_i\})\dd x
 =-2\Re\sum_{k=1}^3\bigg\{&
  \int_0^1\frac{W_{z_{ik}}(x)W_{z_{jk}}(x)
                 -W_{z_{ik}}(0)W_{z_{jk}}(0)}x\dd x\notag\\
 &+\int_1^\infty\frac{W_{z_{ik}}(x)W_{z_{jk}}(x)}x\dd x
   -dW_{z_{ik}}(0)W_{z_{jk}}(0)\bigg\}.
\end{align}
\end{corollary}
All integrals on the right are ordinary convergent integrals. For example,
at $(a_1,a_2,a_3)=(0,1/4,2/3)$ the common value is approximately
\[
 30.589007126644942529054519724157914\ldots.
\]
This decimal is a diagnostic, not a claim of a reduction to a smaller
constant basis. The accompanying code evaluates the left side by its
three local endpoint subtractions and the right side by convergent
polylogarithmic series; their residual at 40 working digits is below
$3\cdot10^{-38}$.

\subsection{Rational colors make the endpoint data finite}
Let $z=e(p/q)\ne1$ with integers $0<p<q$. Splitting the polylogarithm
series into residue classes gives the entire identity
\begin{equation}\label{tsc:finite-DFT}
 \Li_u(z)=q^{-u}\sum_{r=1}^q z^r\zeta(u,r/q).
\end{equation}
The apparent pole on the right cancels because $\sum_{r=1}^qz^r=0$.
Lerch's formula $\zeta^{[1]}(0,x)=\log\Gamma(x)-\ell/2$ then gives
\begin{equation}\label{tsc:W-boundary}
 W_z(0)=\sum_{r=1}^qz^r\log\Gamma(r/q)
       -(\alpha+\log q)\frac{z}{1-z}.
\end{equation}
More generally, substituting \eqref{tsc:finite-DFT} into \eqref{tsc:Vdef}
gives every $V_m(z,0)$ as a finite combination of rational-argument
Hurwitz jets at zero. Thus the endpoint terms in the rational-shift
three-digamma formula are explicitly Gamma-valued, even when the remaining
convergent integrals have not been reduced further.

% END sections/05-mellin-jets

% BEGIN sections/06-primitives.tex
\section{Mean-zero primitives and harmonic coefficients}\label{tsc:primitives}
For $r\ge1$, let $U_{m,r}$ be the unique mean-zero periodic distribution
satisfying $\partial_x^r U_{m,r}=C_m$. Equivalently,
\[
 \widehat U_{m,r}(n)=(2\pi\ii n)^{-r}\widehat C_m(n),\quad n\ne0,
 \qquad \widehat U_{m,r}(0)=0.
\]
These are represented by ordinary locally integrable functions. We give
the exact normalization because a primitive fixed at a base point is not
in general the same function.

For $N\ge0$, define complete homogeneous harmonic coefficients by
\begin{equation}\label{tsc:homogeneous}
 \sum_{j\ge0}h_j^{(N)}v^j
     =\prod_{k=1}^N(1-v/k)^{-1}
     =\exp\left(\sum_{k\ge1}\frac{H_N^{(k)}}k v^k\right).
\end{equation}
Thus $h_0^{(N)}=1$, $h_1^{(N)}=H_N$,
$h_2^{(N)}=(H_N^2+H_N^{(2)})/2$, and
\[
 h_3^{(N)}=\frac{H_N^3+3H_NH_N^{(2)}+2H_N^{(3)}}6.
\]
For $N=0$, all positive-index coefficients are zero.

\begin{proposition}[Explicit normalized primitive]\label{tsc:primitive-formula}
For $m\ge0$, $r\ge1$, and $0<x<1$,
\begin{align}\label{tsc:Ucoef}
 U_{m,r}(x)
 &=m![u^{m+1}]\frac{\Gamma(1+u)}{\Gamma(r+u)}\zeta(1-r-u,x)\notag\\
 &=\frac{(-1)^{m+1}m!}{(r-1)!}
       \sum_{j=0}^{m+1}\frac{h_{m+1-j}^{(r-1)}}{j!}
                   \zeta^{[j]}(1-r,x).
\end{align}
In particular,
\begin{equation}\label{tsc:Uone}
 U_{m,1}(x)=\frac{(-1)^{m+1}}{m+1}\zeta^{[m+1]}(0,x),
 \qquad U_{0,1}(x)=\frac\ell2-\log\Gamma(x).
\end{equation}
\end{proposition}
\begin{proof}
The meromorphic family $\Gamma(u)\calB_{r+u}$ has Fourier coefficients
$\Gamma(u)(2\pi\ii n)^{-r-u}$. Its coefficient $m![u^m]$ is therefore
the stated mean-zero primitive. On the open interval, use
\eqref{tsc:Bhurwitz} and move the factor $1/u$ to the coefficient index;
this proves the first equality in \eqref{tsc:Ucoef}. Next,
\[
 \frac{\Gamma(1+u)}{\Gamma(r+u)}
   =\frac1{(r-1)!}\prod_{k=1}^{r-1}(1+u/k)^{-1}
   =\frac1{(r-1)!}\sum_{j\ge0}(-1)^jh_j^{(r-1)}u^j.
\]
Multiplication by the Taylor expansion of $\zeta(1-r-u,x)$ gives the second
equality. The case $r=1$ and Lerch's formula give \eqref{tsc:Uone}.
Mean zero and the distributional derivative relation follow from the Fourier
construction; they are not extra constants to be chosen afterward.
\end{proof}
These one-function ladders belong to the classical Hurwitz primitive
calculus and its repository developments \cite{tsc:em2,tsc:shj}. Their
three-factor use is the next result.

\begin{theorem}[All primitive triple identities]\label{tsc:primitive-triple}
For $r_i\ge1$ and $m_i\ge0$,
\begin{align}\label{tsc:Utriple}
 \int_0^1\prod_{i=1}^3U_{m_i,r_i}(\{x-a_i\})\dd x
 ={}&(2\pi)^{-r_1-r_2-r_3}\left(\prod_{i=1}^3m_i!\right)\notag\\
 &\quad\cdot[\mathbf u^{\mathbf m+\mathbf1}]
     \left(\prod_{i=1}^3E(u_i)\right)
                         \calT(\mathbf r+\mathbf u;\mathbf a).
\end{align}
This is an ordinary convergent integral. It remains valid when some or all
shifts coincide, with the right side evaluated in its locally convergent
positive-order Tornheim representation.
\end{theorem}
\begin{proof}
For separated shifts, apply \eqref{tsc:master-eq} to
$\prod_i\Gamma(u_i)\calB_{r_i+u_i}$ and extract the three nonnegative
Laurent coefficients, exactly as in Theorem~\ref{tsc:stieltjes-triple}.
This gives \eqref{tsc:Utriple}.

Near a seam, $U_{m,1}$ has at worst logarithmic growth of finite degree.
For $r\ge2$ the possible singular term is a power $x^{r-1}$ multiplied
by a logarithmic polynomial, and the function remains bounded. Every
finite product in the theorem is therefore integrable, including at a
common seam. On the double-sum side, a neighborhood of $r_i\ge1$ is an
absolute-convergence domain even for trivial colors: logarithmic
spectral derivatives are dominated by a slightly decreased positive power,
and the worst central denominator is $mn(m+n)$. Summing first over
$m+n=N$ gives $2H_{N-1}/N^2$. This proves local normal convergence for
all the finitely many derivatives in question. Hence the separated formula
passes to coincident shifts on both sides. This uses no continuation of a
singular Stieltjes product to a collision.
\end{proof}

\begin{corollary}[Shifted centered log-Gamma cube]\label{tsc:loggamma-triple}
Let $G(x)=\log\Gamma(x)-\ell/2$. Then
\begin{equation}\label{tsc:Gtriple}
 \int_0^1\prod_{i=1}^3G(\{x-a_i\})\dd x
   =-\frac1{(2\pi)^3}
       \left.\mathscr D_A\calT(\mathbf s;\mathbf a)\right|_{\mathbf s=(1,1,1)}.
\end{equation}
\end{corollary}
At coincident shifts this is another representation of the known
Bailey--Borwein--Borwein cubic log-Gamma evaluation, not a new evaluation
of that old special case \cite{tsc:bbb,tsc:gamma-chapter}. Arbitrary shifted
mixed-index and mixed-primitive versions follow directly from
\eqref{tsc:Utriple} without another Fourier calculation.

% END sections/06-primitives

% BEGIN sections/07-contacts.tex
\section{Spatial derivatives and exact harmonic contact terms}\label{tsc:contacts}
The identity $\partial_x\calB_s=\calB_{s-1}$ is distributional. Passing
from it to a pointwise polygamma product requires a local correction. This
section gives that correction explicitly and does not alter the correct
pointwise derivative tower already present in the repository.

\subsection{Finite reductions at negative integral spectral orders}
For $r\ge1$, $\calB_{-r}=\delta_0^{(r)}$. Applying this to the master
identity gives a useful family of exact spectral resonances.
\begin{proposition}\label{tsc:negative-resonance}
For $r\ge1$ and any complex $s,t$,
\begin{align}\label{tsc:negative-resonance-eq}
 \calT(-r,s,t;\mathbf a)
 ={}&(2\pi)^{s+t-r}(-1)^r
   \sum_{j=0}^r\binom rj
    \calB_{s-j}(a_1-a_2)\calB_{t-r+j}(a_1-a_3).
\end{align}
Here the functions on the right are evaluated away from their seam using
\eqref{tsc:Bhurwitz}. At order zero, the missing constant mode changes
the formula to
\begin{align}\label{tsc:zero-resonance}
 \calT(0,s,t;\mathbf a)={}&(2\pi)^{s+t}
    \calB_s(a_1-a_2)\calB_t(a_1-a_3)\notag\\
 &-e^{-\pi\ii(s-t)/2}\Li_{s+t}(e(a_3-a_2))\notag\\
 &-e^{\pi\ii(s-t)/2}\Li_{s+t}(e(a_2-a_3)).
\end{align}
\end{proposition}
\begin{proof}
Pair $\delta_{a_1}^{(r)}$ with the product of the other two, smooth,
factors. The sign $(-1)^r$ is the defining sign of a delta derivative,
and the product rule with $\partial_x^j\calB_s=\calB_{s-j}$ gives
\eqref{tsc:negative-resonance-eq}. For zero order use
$\calB_0=\delta_0-1$, not $\delta_0$. The remaining bilinear mean is
computed by its two opposite-frequency sectors and equals the two
polylogarithm terms after multiplication by $(2\pi)^{s+t}$.
\end{proof}
The formulas reduce the complete six-sector expression on these spectral
hyperplanes. They do not determine its transverse derivatives, which are
precisely among the data needed for general Stieltjes triple products.

\subsection{Elementary harmonic coefficients}
Define $e_j^{(r)}$ by
\begin{equation}\label{tsc:elementary}
 \prod_{k=1}^r(1+v/k)=\sum_{j=0}^r e_j^{(r)}v^j,
 \qquad e_j^{(r)}=0\quad(j>r).
\end{equation}
Thus $e_1^{(r)}=H_r$,
$e_2^{(r)}=(H_r^2-H_r^{(2)})/2$, and
$e_3^{(r)}=(H_r^3-3H_rH_r^{(2)}+2H_r^{(3)})/6$.
Unlike the complete coefficients in the primitive ladder, these have finite
support in $j$.

Let $D_{m,r}$ denote the scale-one, one-sided Hadamard extension of the
\emph{pointwise} function $\gamma_m^{(r)}(x)$ on $0<x<1$. At the right-hand
seam, expand the test function through degree $r$, integrate all singular
power-log terms, and delete the divergent terms in the cutoff. This
specifies the extension without arbitrary delta terms. Put $D_{m,0}=C_m$.

\begin{theorem}[Derivative--extension conversion]\label{tsc:contact-single}
For $r\ge0$,
\begin{equation}\label{tsc:single-contact}
 \partial_x^r C_m=D_{m,r}+c_{m,r}\delta_0^{(r)},
 \qquad c_{m,r}=(-1)^{m+1}m!e_{m+1}^{(r)}.
\end{equation}
In particular $c_{m,0}=0$, $c_{0,r}=-H_r$, and $c_{m,r}=0$ for $m\ge r$.
\end{theorem}
\begin{proof}
In a local right-hand coordinate, the meromorphic distribution
$x_+^{u-r-1}$ has residue
$(-1)^r\delta_0^{(r)}/r!$ at $u=0$. Its regular Laurent coefficients
are the scale-one finite parts of $x^{-r-1}\log^j x/j!$.
Differentiating the singular term $x^{u-1}$ in
$\zeta(1-u,x)$ gives
\[
 \partial_x^r x^{u-1}=(-1)^r r!P_r(u)x^{u-r-1},
 \qquad P_r(u)=\prod_{k=1}^r(1-u/k).
\]
Hence the residue contribution after multiplication is
$P_r(u)\delta_0^{(r)}/u$. The regular nonresidue terms give exactly the
raw extensions $D_{m,r}$. Comparing with the $r$th derivative of
\eqref{tsc:Zlaurent}, the additional coefficient is
$m![u^{m+1}]P_r(u)=(-1)^{m+1}m!e_{m+1}^{(r)}$.
The remaining smooth local terms introduce no contact distribution.
For $r=0$ the assertion is immediate. This proves the formula.
\end{proof}

\begin{theorem}[Contact-complete derivative triple]\label{tsc:contact-triple}
Let $r_i,m_i\ge0$, $R=r_1+r_2+r_3$, and retain distinct shifts. Then
\begin{align}\label{tsc:contact-triple-eq}
 &\FP\int_0^1\prod_{i=1}^3
            \gamma_{m_i}^{(r_i)}(\{x-a_i\})\dd x\notag\\
 &\quad=(2\pi)^R\left(\prod_{i=1}^3m_i!\right)
    [\mathbf u^{\mathbf m+\mathbf1}]
        \left(\prod_{i=1}^3E(u_i)\right)
                       \calT(\mathbf u-\mathbf r;\mathbf a)\notag\\
 &\qquad\quad-
    \sum_{i=1}^3(-1)^{r_i}c_{m_i,r_i}
       \left.\partial_x^{r_i}
              \prod_{j\ne i}\gamma_{m_j}^{(r_j)}(\{x-a_j\})
       \right|_{x=a_i}.
\end{align}
The local derivatives in the last line are ordinary derivatives of smooth
functions at $a_i$. In particular, setting all $m_i=0$ gives identities for
arbitrary products of three separated polygamma functions.
\end{theorem}
\begin{proof}
By \eqref{tsc:single-contact}, each raw factor is
$D_{m_i,r_i,a_i}=\partial_x^{r_i}C_{m_i,a_i}
                 -c_{m_i,r_i}\delta_{a_i}^{(r_i)}$.
In the product expansion, every term containing two delta-supported
factors is zero because their supports are disjoint. A term containing
one delta derivative pairs with the smooth remaining product, producing
$(-1)^{r_i}$ times its ordinary $r_i$th derivative at $a_i$. These are
exactly the terms in the final line of \eqref{tsc:contact-triple-eq}.

The product of the three periodic derivatives is
$(-1)^R\partial_{a_1}^{r_1}\partial_{a_2}^{r_2}\partial_{a_3}^{r_3}
\calJ_{\mathbf m}$. Apply \eqref{tsc:lowering} to
\eqref{tsc:all-index}. The factor $(-1)^R$ cancels the sign of
$(-2\pi)^R$, giving the first line of the right-hand side.
\end{proof}

The correction is not optional. For instance,
$\partial_x C_0=D_{0,1}-\delta_0'$. Omitting it already changes a triple
with one spatial derivative by the derivative of the product of the two
other endpoint values. The theorem explains the correction without
claiming that the manuscript's pointwise formulas are wrong.

For explicit simplification of endpoint derivatives, the classical
Pochhammer identity yields, for $r\ge1$,
\begin{equation}\label{tsc:pointwise-tower}
 \gamma_n^{(r)}(x)=(-1)^{r+n}r!n!
      \sum_{k=0}^{\min(r,n)}
       e_k^{(r)}\frac{\zeta^{[n-k]}(r+1,x)}{(n-k)!}.
\end{equation}
Indeed, expand
$\partial_x^r\zeta(1+v,x)=(-1)^r(1+v)_r\zeta(1+r+v,x)$ at zero.
This derivation also fixes every factorial and sign in
\eqref{tsc:pointwise-tower}.

% END sections/07-contacts

% BEGIN sections/08-reflection.tex
\section{Reflected polygamma products at arbitrary depth}\label{tsc:reflected}
The generic triple has a double-sum coordinate class. A special reflected
subalgebra reduces much further. Its basic kernel is
\begin{equation}\label{tsc:Kdef}
 K(x)=\PV\,\pi\cot(\pi x)=C_0(x)-C_0(-x).
\end{equation}
Its zero Fourier coefficient is zero, and its other coefficients are
$-\pi\ii\sgn n$. The reflection formula gives the pointwise identity
$\psi(1-x)-\psi(x)=\pi\cot\pi x$. The two one-sided logarithmic
counterterms cancel in the reflected difference, giving the symmetric
principal value in \eqref{tsc:Kdef}.

\subsection{An all-index Stieltjes--cotangent transform}
For $0<a<1$, put
\begin{equation}\label{tsc:Idef}
 I_m(a)=\langle C_m(x)K(x-a),1\rangle.
\end{equation}
This means the one-sided finite part at zero and the symmetric principal
value at $a$; these singularities are separated.

\begin{proposition}[Spectral generating function]\label{tsc:Igen}
Near $u=0$,
\begin{align}\label{tsc:Igen-eq}
 \sum_{m\ge0}\frac{I_m(a)}{m!}u^m
  =\pi\csc(\pi u)\zeta(1-u,1-a)
     -\pi\cot(\pi u)\zeta(1-u,a)+\frac{K(a)}u.
\end{align}
The singularity on the right is removable only after the terms are combined.
\end{proposition}
\begin{proof}
For $\Re u>1$, Fourier pairing gives
\begin{align*}
 \langle\Gamma(u)\calB_u(x)K(x-a),1\rangle
 =\frac{\pi\ii\Gamma(u)}{(2\pi)^u}
   \left[e^{-\pi\ii u/2}\Li_u(e(a))
               -e^{\pi\ii u/2}\Li_u(e(-a))\right].
\end{align*}
The two Hurwitz Fourier formulas at $a$ and $1-a$ transform this into
the first two terms of \eqref{tsc:Igen-eq}. Analytic continuation in the
separated product extends the identity to zero. The residue of
\eqref{tsc:Zlaurent} paired with $K(x-a)$ is
$K(-a)-\langle K,1\rangle=-K(a)$. Adding $K(a)/u$ removes precisely that
residue and leaves the generating series for $I_m$.
\end{proof}
This is a bilinear kernel identity, compatible with the repository's
existing closure machinery; it is used as an input to the all-depth
reduction rather than advertised as a new bilinear principle.

\begin{corollary}[Finite single-Stieltjes formula]\label{tsc:Iexplicit}
For all $m\ge0$,
\begin{align}\label{tsc:I-formula}
 I_m(a)={}&\frac{\gamma_{m+1}(1-a)-\gamma_{m+1}(a)}{m+1}\notag\\
 &+2m!\sum_{k=1}^{\lfloor(m+1)/2\rfloor}
       \frac{\zeta(2k)}{(m-2k+1)!}
       \left[(1-2^{1-2k})\gamma_{m-2k+1}(1-a)
                       +\gamma_{m-2k+1}(a)\right]\notag\\
 &-\mathbf1_{2\mid m}\,2m!(2-2^{-m-1})\zeta(m+2).
\end{align}
\end{corollary}
\begin{proof}
Insert \eqref{tsc:stieltjes-convention} into \eqref{tsc:Igen-eq} and use
\[
 \pi\csc(\pi u)=\frac1u+
     2\sum_{k\ge1}(1-2^{1-2k})\zeta(2k)u^{2k-1},\qquad
 \pi\cot(\pi u)=\frac1u-2\sum_{k\ge1}\zeta(2k)u^{2k-1}.
\]
The remaining pole cancels by
$\gamma_0(1-a)-\gamma_0(a)=-K(a)$. The products with the spectral pole
$-1/u$ contribute the last line of \eqref{tsc:I-formula}; retaining them
is necessary even for $m=0$. Coefficient comparison proves the formula.
\end{proof}
In low degrees,
\begin{align}\label{tsc:Iexamples}
 I_0(a)&=\gamma_1(1-a)-\gamma_1(a)-\frac{\pi^2}{2},\notag\\
 I_1(a)&=\frac{\gamma_2(1-a)-\gamma_2(a)}2
           +\zeta(2)\bigl(\gamma_0(1-a)+2\gamma_0(a)\bigr),\notag\\
 I_2(a)&=\frac{\gamma_3(1-a)-\gamma_3(a)}3
       +2\zeta(2)\gamma_1(1-a)+4\zeta(2)\gamma_1(a)
       -\frac{15}{2}\zeta(4).
\end{align}
For example, $I_0(1/2)=-\pi^2/2$ and
$I_1(1/2)=\pi^2(\gamma+2\log2)/2$. The machine-readable table in
\src{results/stieltjes_cotangent_coefficients.json} extends through $m=12$.

\subsection{A partial-fraction identity at every number of factors}
Let $b_1,\ldots,b_N$ be pairwise distinct on the circle and define
\begin{equation}\label{tsc:Acot}
 A_j(\mathbf b)=\prod_{\ell\ne j}K(b_j-b_\ell).
\end{equation}
Empty products equal one.
\begin{lemma}[All-depth cotangent partial fractions]\label{tsc:cot-PF}
As a symmetric finite-part distribution,
\begin{equation}\label{tsc:cot-PF-eq}
 \prod_{j=1}^N K(x-b_j)
     =\pi^N\cos\frac{N\pi}{2}
        +\sum_{j=1}^NA_j(\mathbf b)K(x-b_j).
\end{equation}
Moreover,
\begin{equation}\label{tsc:cot-residue-sum}
 \sum_{j=1}^NA_j(\mathbf b)=\pi^{N-1}\sin\frac{N\pi}{2}.
\end{equation}
\end{lemma}
\begin{proof}
Put $X=e(x)$ and $B_j=e(b_j)$. Then
$K(x-b_j)=\pi\ii(X+B_j)/(X-B_j)$.
The product is a rational function with only simple poles at $B_j$;
subtracting $A_jK(x-b_j)$ removes its principal part at that pole.
The difference has no finite pole and is bounded at infinity, hence is a
constant. As $X\to0$ the kernels tend to $-\pi\ii$, and as $X\to\infty$
they tend to $\pi\ii$. Adding and subtracting the two limiting equations
gives the constant and the sum in \eqref{tsc:cot-PF-eq}--\eqref{tsc:cot-residue-sum}.
Because the principal parts have matched at every real pole, the rational
identity passes to the specified symmetric extensions without a delta
correction.
\end{proof}

\begin{theorem}[Finite reduction at arbitrary reflected depth]\label{tsc:all-depth}
Assume $a,b_1,\ldots,b_N$ are pairwise distinct modulo one. Then
\begin{equation}\label{tsc:all-depth-eq}
 \boxed{\displaystyle
 \FP\int_0^1 C_m(x-a)\prod_{j=1}^NK(x-b_j)\dd x
       =\sum_{j=1}^NA_j(\mathbf b)I_m(\{b_j-a\}).}
\end{equation}
In addition,
\begin{equation}\label{tsc:cot-mean}
 \PV\int_0^1\prod_{j=1}^NK(x-b_j)\dd x
                 =\pi^N\cos\frac{N\pi}{2}.
\end{equation}
In \eqref{tsc:all-depth-eq}, $C_m$ denotes its prescribed extension, so the
integrand away from its seam is the ordinary function $\gamma_m$.
\end{theorem}
\begin{proof}
Multiply \eqref{tsc:cot-PF-eq} by $C_m(x-a)$ and pair with one. The
constant term vanishes because $C_m$ has mean zero. Translate $x$ by $a$
to identify each remaining pairing with $I_m(\{b_j-a\})$. Taking the
mean without $C_m$ proves \eqref{tsc:cot-mean}.
\end{proof}

\subsection{Reflected polygamma derivatives}
The balanced pointwise combination is
\begin{equation}\label{tsc:reflected-poly}
 K^{(r)}(x)=(-1)^r\psi^{(r)}(1-x)-\psi^{(r)}(x).
\end{equation}
Its symmetric extension is exactly the distributional derivative of $K$:
differentiation of the symmetric principal part $1/x$ gives the symmetric
Hadamard principal parts, without an additional delta term. Consequently,
for arbitrary $r_j\ge0$ the integral obtained by replacing every
$K(x-b_j)$ in \eqref{tsc:all-depth-eq} with $K^{(r_j)}(x-b_j)$ equals
\begin{equation}\label{tsc:reflected-derivatives}
 \prod_{j=1}^N(-\partial_{b_j})^{r_j}
           \left(\sum_{j=1}^NA_j(\mathbf b)I_m(\{b_j-a\})\right).
\end{equation}
The equality follows by differentiating a product of separated distributions;
it is not an unqualified interchange with an unregularized improper integral.
Formula \eqref{tsc:I-formula}, the product rule, and
\eqref{tsc:pointwise-tower} make the result finite in Hurwitz-zeta jets and
trigonometric coefficients. This provides explicit identities for arbitrary
numbers and orders of reflected polygamma factors.

% END sections/08-reflection

% BEGIN sections/09-gamma.tex
\section{Rational Gamma reductions and explicit antiderivatives}\label{tsc:gamma}
At Stieltjes index zero, the antisymmetric first Stieltjes constant has a
finite Gamma reduction. We include a proof to make the all-depth evaluation
self-contained.

\begin{proposition}[Rational antisymmetric first Stieltjes value]\label{tsc:rational-D}
For integers $0<p<q$, with $a=p/q$,
\begin{align}\label{tsc:rational-D-eq}
 \gamma_1(1-a)-\gamma_1(a)
 ={}&\pi\cot(\pi a)\bigl(\gamma+\log(2\pi q)\bigr)\notag\\
 &-2\pi\sum_{r=1}^{q-1}\sin(2\pi pr/q)\log\Gamma(r/q).
\end{align}
Coprimality of $p,q$ is not required.
\end{proposition}
\begin{proof}
Write $z=e(a)$ and
$P(u)=\Gamma(u)(2\pi)^{-u}$. Subtract the two Hurwitz Fourier expansions:
\[
 \zeta(1-u,1-a)-\zeta(1-u,a)
      =2\ii P(u)\sin(\pi u/2)
                    \bigl(\Li_u(z)-\Li_u(z^{-1})\bigr).
\]
Now $2\ii P(u)\sin(\pi u/2)=\pi\ii(1-Au+O(u^2))$.
For real $u$, the difference of polylogarithms is $2\ii\Im\Li_u(z)$.
Comparing coefficients of $u$ gives
\[
 \gamma_1(1-a)-\gamma_1(a)
       =\pi A\cot(\pi a)-2\pi\Im\Li_0^{[1]}(z),
\]
since $\Im\Li_0(z)=\tfrac12\cot(\pi a)$.
Differentiate \eqref{tsc:finite-DFT} at zero. The constant part of Lerch's
formula vanishes in the finite sum, so
\[
 \Im\Li_0^{[1]}(z)=
     \sum_{r=1}^{q-1}\sin(2\pi pr/q)\log\Gamma(r/q)
                -\frac{\log q}{2}\cot(\pi a).
\]
Substitution proves \eqref{tsc:rational-D-eq}.
\end{proof}
This rational reflection formula is a consequence of classical Hurwitz
Fourier and Gamma identities. Its use here is to finish a whole family of
reflected product evaluations, rather than to claim a new isolated formula
for $\gamma_1(p/q)$.

\begin{theorem}[All-depth rational Gamma evaluation]\label{tsc:rational}
Let $a,b_1,\ldots,b_N$ be pairwise distinct rational circle points. Choose
$q$ with $d_j=\{b_j-a\}=p_j/q$, $0<p_j<q$. Then
\begin{align}\label{tsc:all-depth-gamma}
 &\FP\int_0^1 C_0(x-a)\prod_{j=1}^NK(x-b_j)\dd x\notag\\
 &\quad=\sum_{j=1}^N A_j(\mathbf b)
   \left\{K(d_j)\bigl(\gamma+\log(2\pi q)\bigr)
      -2\pi\sum_{r=1}^{q-1}\sin(2\pi p_jr/q)\log\Gamma(r/q)
      -\frac{\pi^2}{2}\right\}.
\end{align}
Thus this arbitrary-depth subclass has a finite expression in
$\gamma,\pi$, logarithms, and log-Gamma values at rational arguments,
with algebraic trigonometric coefficients.
\end{theorem}
\begin{proof}
Insert \eqref{tsc:rational-D-eq} into $I_0(d)$ in
\eqref{tsc:Iexamples}, then into Theorem~\ref{tsc:all-depth}.
At rational angles, the trigonometric coefficients belong to cyclotomic
number fields. No arithmetic independence statement is used.
\end{proof}
For even $N$, the final constants sum to zero by
\eqref{tsc:cot-residue-sum}. For odd $N$, their sum is
$-\pi^{N+1}\sin(N\pi/2)/2$. The parity effect checks the normalization.

\subsection{A cubic identity with only interior principal values}
\begin{theorem}[Quarter--half cubic digamma integral]\label{tsc:explicit-cubic}
The following principal value exists and satisfies
\begin{align}\label{tsc:explicit-cubic-eq}
 \boxed{\displaystyle
 \PV\int_0^1\psi(x)
     \cot\!\bigl(\pi(x-\tfrac14)\bigr)
     \cot\!\bigl(\pi(x-\tfrac12)\bigr)\dd x
  =\gamma+4\log2+3\log\pi-4\log\Gamma(\tfrac14).}
\end{align}
The principal values are symmetric at $1/4$ and $1/2$. The singularity at
zero is removable; the integrand tends to $-\pi$ there.
\end{theorem}
\begin{proof}
Here $A_1=K(-1/4)=-\pi$ and $A_2=K(1/4)=\pi$. The partial fraction
identity is
\[
 K(x-\tfrac14)K(x-\tfrac12)
        =-\pi^2-\pi K(x-\tfrac14)+\pi K(x-\tfrac12).
\]
Since $C_0=-\psi$, Theorem~\ref{tsc:all-depth} shows that the left side of
\eqref{tsc:explicit-cubic-eq} is
\[
 \frac{I_0(1/4)-I_0(1/2)}\pi
      =\frac{\gamma_1(3/4)-\gamma_1(1/4)}\pi.
\]
At $p/q=1/4$, Proposition~\ref{tsc:rational-D} gives
\[
 \frac{\gamma_1(3/4)-\gamma_1(1/4)}\pi
  =\gamma+\log(8\pi)-2\log\Gamma(1/4)+2\log\Gamma(3/4).
\]
Euler reflection, $\Gamma(1/4)\Gamma(3/4)=\pi\sqrt2$, reduces this to
the claimed right side.

Finally, $\psi(x)=-1/x+O(1)$,
$\cot(\pi(x-1/2))=-\pi x+O(x^3)$, and
$\cot(\pi(x-1/4))=-1+O(x)$. Their product tends to $-\pi$.
Consequently the one-sided finite-part notation needed in the general
theorem disappears in this specialization: only the two stated interior
principal values remain.
\end{proof}
The common value is
\[
 1.63190394589720479122328086909541565458453405168\ldots.
\]
The evaluation is proved by the preceding identities; the decimal is a
separate quadrature check.

\subsection{A log-Gamma cotangent transform}
\begin{proposition}[An ordinary primitive of the reflected kernel]\label{tsc:Gcot}
For $0<a<1$,
\begin{align}\label{tsc:Gcot-eq}
 \PV\int_0^1\log\Gamma(x)K(x-a)\dd x
   =\frac{\zeta^{[2]}(0,a)+\zeta^{[2]}(0,1-a)}2
        +\frac{\pi^2}{2}\left(a-\frac12\right).
\end{align}
The endpoint at zero is ordinarily integrable. In particular,
\begin{equation}\label{tsc:Gtan}
 \PV\int_0^1\log\Gamma(x)\tan(\pi x)\dd x
    =\frac{\log2\log(2\pi)+\tfrac12\log^22}{\pi}.
\end{equation}
\end{proposition}
\begin{proof}
Let the left side of \eqref{tsc:Gcot-eq} be $L(a)$.
By \eqref{tsc:Uone}, the periodic distributional derivative of
$\log\Gamma(x)-\ell/2$ is $-C_0$. Differentiation of the separated
pairing, followed by distributional integration by parts, gives
$L'(a)=-I_0(a)$.
Differentiating $\partial_x\zeta(s,x)=-s\zeta(s+1,x)$ at $s=0$, with
the spectral pole already canceled, gives
\[
 \partial_x\zeta^{[2]}(0,x)=2\gamma_1(x).
\]
Hence the derivative of the proposed right side is
$\gamma_1(a)-\gamma_1(1-a)+\pi^2/2=-I_0(a)$.
Both sides have mean zero in $a$: this is the zero mean of $K$ on the
left and of the Hurwitz jets at zero on the right. Their constant difference
is therefore zero.

At $a=1/2$, use $K(x-1/2)=-\pi\tan\pi x$ and
$\zeta(s,1/2)=(2^s-1)\zeta(s)$. Twice differentiating at zero yields
$\zeta^{[2]}(0,1/2)=-\log2\log(2\pi)-\tfrac12\log^22$, which proves
\eqref{tsc:Gtan}.
\end{proof}

More generally, an explicit pointwise antiderivative of every $I_m(a)$ in
\eqref{tsc:I-formula} is obtained by replacing $\gamma_n(a)$ with
\[
 Q_n(a)=\frac{(-1)^{n+1}}{n+1}\zeta^{[n+1]}(0,a),\qquad Q_n'(a)=\gamma_n(a),
\]
replacing $\gamma_n(1-a)$ with $-Q_n(1-a)$, and integrating the constant
term linearly. This produces a finite Hurwitz-jet antiderivative at every
Stieltjes index. A chosen base-point constant or a mean-zero periodic
normalization must then be specified separately.

% END sections/09-gamma

% BEGIN sections/10-audit.tex
\section{Audit findings, claim boundaries, and reproducibility}\label{tsc:audit}
\subsection{A proposed update, not a fabricated erratum}
The inspected shifted-jet report explicitly asks for the coordinate class
of separated three-factor Stieltjes products. Theorems~\ref{tsc:stieltjes-triple}
and \ref{tsc:mellin-all} answer that question affirmatively, and
\eqref{tsc:digamma-operator} and \eqref{tsc:psi-integrals} supply its first
three-digamma target. A suitable manuscript update is therefore to mark the
\emph{representation question} resolved and replace it with the sharper
questions about collisions and further period reductions in the next
section.

No false theorem was identified in the selected canonical passages read
for this continuation. In particular, their pointwise derivative tower is
correct, the existing cubic log-Gamma attribution is correct, and the
current $S_4$ status must remain proved. The caveats below are mathematical
integration safeguards, not claims that the repository currently contains
all of the corresponding mistakes.

\paragraph{Preserve the constant mode.}
The normalized zero-order family is $\calB_0=\delta_0-1$. Replacing it by
$\delta_0$ would give a false value for the triple kernel. The exact test
$\calT(0,0,0)=2$ detects this normalization error.

\paragraph{Preserve the endpoint constants.}
In \eqref{tsc:I-formula}, the even-index zeta term comes from multiplying a
spectral pole by a regular trigonometric coefficient. In
\eqref{tsc:Mfunctional}, the finite constant comes from the Mellin pole
$F(0)/u$. Omitting either term changes the identity; in the first case,
$I_0(1/2)=-\pi^2/2$ is an immediate counterexample.

\paragraph{Distinguish discrete and analytic order formulas.}
The finite reduction \eqref{tsc:integer-reduction} has integer parameters
$r,s$. It cannot simply be differentiated in those parameters. The entire
colored family and its subtracted Mellin integrals provide the required
analytic replacement. Values on isolated integer hyperplanes do not fix
transverse spectral derivatives.

\paragraph{Do not transport pointwise derivatives silently.}
A derivative of a periodic finite-part extension and a raw finite part of
the pointwise derivative differ by \eqref{tsc:single-contact}. The triple
correction is \eqref{tsc:contact-triple-eq}. Conversely, symmetric derivatives
of the reflected cotangent kernel have their own compatible convention.
They should not be assigned the one-sided correction of an unreflected
Stieltjes factor.

\paragraph{Keep collision claims separate.}
The generic Stieltjes product requires pairwise distinct shifts. Setting
colors equal to one in the global entireness theorem is invalid. Ordinary
primitive products can allow coincident shifts for the different,
explicitly proved convergence reason in Theorem~\ref{tsc:primitive-triple}.

\subsection{What the tests establish}
The exact checker uses Gaussian rational arithmetic for fourth-root Fourier
phases and SymPy exact arithmetic for polynomial and harmonic coefficients.
Its 621 assertions comprise the following finite components.
\begin{center}
\begin{tabular}{lr}
\toprule
Exact component & Assertions\\\midrule
Six-sector finite Fourier identities & 128\\
Deliberately corrupted-formula rejection controls & 3\\
Zero-order six-sector value $2$ & 60\\
Integer Tornheim partial-fraction identities & 64\\
Cotangent partial fractions and residue sums & 20\\
Stieltjes--cotangent coefficients, $0\le m\le12$ & 13\\
Complete homogeneous harmonic coefficients & 143\\
Elementary harmonic coefficients & 91\\
Normalized primitive derivative ladder & 99\\\midrule
Total & 621\\\bottomrule
\end{tabular}
\end{center}
These checks do not constitute a formal proof of distributional analytic
continuation. That proof is the argument in the article. The test scope is
recorded explicitly in \src{results/exact_checks.json}.

The numerical suite makes 56 independent comparisons. It includes direct
subtracted quadrature of $I_0$ and the log-Gamma transform, rational Gamma
reductions, products with two through four cotangent factors, and a
nonzero-parameter check of the Hurwitz generating identity. For three
positive integral spectral triples, the six Mellin sectors are compared
with exactly integrated Bernoulli polynomials. The rational targets are
\begin{equation}\label{tsc:bernoulli-targets}
 \frac8{10125},\qquad -\frac{11}{15360},\qquad
 \frac{6119}{63026250};
\end{equation}
the associated orders and shifts are recorded in the JSON results.
Finally, three singular digamma triples are evaluated both by local
endpoint subtractions and by \eqref{tsc:psi-integrals}. The polylogarithm
order derivatives in that check are computed from convergent series,
not by infinitesimal black-box differentiation of a polylogarithm routine.

The default environment used Python 3.13.5, SymPy 1.14.0, and mpmath 1.3.0.
Most numerical checks use 50 decimal working digits; the more expensive
Mellin and polylogarithmic triple checks use 40. Acceptance thresholds range
from 31 to 40 relative decimal digits, depending on the test. None of these
are certified interval calculations. High-index Stieltjes triples and the
full derivative triple family are proved analytically here; they are not
claimed to have been exhaustively numerically evaluated.

\subsection{Reproduction and integration}
From the package directory, run
\begin{verbatim}
python -m pip install -r requirements.txt
python code/verify_exact.py
python code/verify_numerical.py --dps 50
python code/build_article.py
\end{verbatim}
The source is modular, with a flattened standalone TeX copy for convenience.
The archive includes the compiled article, coefficient tables, executed
check logs, source provenance, an additive integration note, and a SHA-256
manifest. No network access is needed by the checking scripts after their
dependencies are installed.

A proposed report destination is
\src{Analysis/Polylogarithms/docs/reports/stieltjes-correlation-calculus/continuations/triple-stieltjes-tornheim/}.
All theorem labels and bibliography keys use the prefix \src{tsc:}.
The integration note distinguishes a proposed status update from an applied
patch. Archive preservation and successful replay should remain separate
from mathematical acceptance into the canonical manuscript.

% END sections/10-audit

% BEGIN sections/11-further.tex
\section{Further research questions}\label{tsc:further}
\subsection{Joint triple collisions and normalization conversion}
For two approaching seams, the separated three-factor formula develops a
local singularity. When all three approach one another, two relative
scales are available. Determine the full two-parameter local expansion,
including the delta derivatives and finite constants needed to compare
separated analytic continuation with a chosen same-point Hadamard product.
A first concrete target is the case $(m_1,m_2,m_3)=(0,0,0)$ along
$a_2-a_1=\lambda\varepsilon$, $a_3-a_1=\varepsilon$ with fixed
$\lambda\notin\{0,1\}$. The dependence on $\lambda$ must be computed,
not erased by an unproved path-independent finite-part prescription.

\subsection{Transverse Tornheim jets and genuine arithmetic reductions}
Proposition~\ref{tsc:negative-resonance} gives finite reductions on integral
spectral hyperplanes, but the Stieltjes formula also uses derivatives
transverse to them. Find additional functional identities that reduce
specific transverse jets at rational colors. The question should be asked
first at a fixed conductor and a fixed total spectral degree. A negative
result must distinguish nonmembership in an explicitly generated formal
relation space from independence of the resulting numerical periods.

\subsection{Four unreflected factors and balanced frequency sectors}
Four nonzero frequencies summing to zero have sectors of types $3+1$ and
$2+2$. The latter are not a direct copy of the Tornheim double sum. Develop
an explicit, jointly regularized coordinate family for these sectors and
a counterpart of Theorem~\ref{tsc:mellin-all}. A useful first target is
four digamma factors at four distinct rational shifts, with every endpoint
constant and contact correction specified. The reflected finite-depth
formula already proved here should serve as a stringent specialization
check, not as evidence that the generic family has depth one.

\subsection{Rational traces that create collisions}
Integer dilation and conductor operations can turn previously distinct
shifts into coincident shifts. Determine the exact trace identities for
the three-factor kernel across that boundary. The existing bilinear
conductor laws and the separated lowering law are useful inputs, but they
do not by themselves determine the new collision counterterms.

\subsection{Certified polylogarithmic jet evaluation}
The present stable series evaluation avoids a known source of numerical
fragility: differentiating a polylogarithm routine at an infinitesimal
order increment. Replace its floating-point truncations with explicit
complex interval enclosures, including a uniform Jonqui\`ere-series
remainder for the unit-color region. Combine that with the subtracted
Mellin integral to obtain rigorous numerical certificates for generic
three-digamma values. This is an algorithmic refinement of proved
identities, not a substitute for their analytic proof.

\subsection{A precise bridge to the remaining fixed-weight candidates}
The repository's $S_6$ and $S_8$ questions concern particular polylogarithmic
period reductions. The present triple calculus does not automatically
settle them. Search for an exact transformation that expresses a candidate's
remaining obstruction in one of the rational-color transverse jets above,
with all boundary terms retained. Such a bridge, together with a proved
jet reduction, would be a meaningful route to a certificate. Numerical
agreement or a shared use of polylogarithms is not such a bridge.

\subsection{Formalization with analytic hypotheses exposed}
A useful formal development would separate the polynomial-growth Fourier
construction, the local meromorphic extension of $x_+^{u-r-1}$, the
separated product lemma, and the finite coefficient algebra. The latter
already has executable exact tests, but not proof-assistant certificates.
The first analytic formal target can be
\eqref{tsc:single-contact}; the trilinear master theorem can then be proved
conditional on the ordinary absolutely convergent Fourier identity before
formalizing the entire continuation.

\section*{Conclusion}
\addcontentsline{toc}{section}{Conclusion}
Three separated Stieltjes singularities have an exact, all-index colored
Tornheim description, and their Laurent coefficients have a convergent
polylogarithmic Mellin representation with explicit endpoint constants.
Mean-zero integration and contact-corrected differentiation extend that
same kernel to Gamma and polygamma products. Reflection selects an
arbitrary-depth subclass with finite single-Stieltjes, and at rational
index-zero shifts finite Gamma, reductions. The resolved coordinate-class
question is therefore a concrete advance within the inspected research
programme, while collisions, transverse arithmetic reductions, and the
remaining fixed-weight conjectures retain clearly stated open targets.

% END sections/11-further

\appendix
% BEGIN sections/12-computation.tex
\section{Independent computation formulas}\label{tsc:computation}
\subsection{Direct three-digamma finite parts}
Suppose $0\le a_1<a_2<a_3<1$, put $a_4=a_1+1$, and let
$L_j=a_{j+1}-a_j$. On the interval beginning at $a_j$, use the coordinate
$t=x-a_j$ and offsets $d_{j\ell}=\{a_j-a_\ell\}$, with $d_{jj}=0$.
Set
\[
 c_j=\prod_{\ell\ne j}\gamma_0(d_{j\ell}).
\]
Then a directly integrable version of the canonical triple is
\begin{equation}\label{tsc:direct-triple}
 \calJ_{000}(\mathbf a)=\sum_{j=1}^3\left\{
   \int_0^{L_j}\left[\prod_{\ell=1}^3\gamma_0(t+d_{j\ell})
                              -\frac{c_j}{t}\right]\dd t
                          +c_j\log L_j\right\}.
\end{equation}
No fractional part is applied inside a single interval: all the offsets
stay in the appropriate left-limit branches until its upper endpoint.
The removable value of the bracket at zero is
\[
 \gamma c_j+
  \sum_{\ell\ne j}\gamma_0'(d_{j\ell})
             \prod_{h\ne j,\ell}\gamma_0(d_{jh}).
\]
The three-digamma product is $-\calJ_{000}$. This quadrature uses no
Tornheim function, Mellin continuation, or polylogarithm derivative, making
it independent of the right side of \eqref{tsc:psi-integrals}.

\subsection{Stable evaluation of the order-zero polylogarithm jet}
For $\mu=-x+\ii\theta$ with $|\mu|<2\pi$, differentiating the classical
Jonqui\`ere expansion \cite[Eq.~25.12.12]{tsc:dlmf-poly} gives
\begin{equation}\label{tsc:Li-order-zero}
 \Li_0^{[1]}(e^\mu)
   =-\frac{\gamma+\Log(-\mu)}{\mu}
       +\sum_{k\ge0}\frac{\zeta^{[1]}(-k)}{k!}\mu^k.
\end{equation}
The principal branch is used for $\Log(-\mu)$.
For $0\le x\le1$ and $0<|\theta|\le\pi$, this convergent disk contains
the entire evaluation region. The coefficients can be generated without
spectral differencing of a polylogarithm:
\begin{align}\label{tsc:negative-zeta-derivative}
 \zeta^{[1]}(0)&=-\frac\ell2,\notag\\
 \frac{\zeta^{[1]}(-k)}{k!}
  &=\frac{(-1)^{k/2}\zeta(k+1)}{2(2\pi)^k},
                                      &&k\ge2\text{ even},\notag\\
 \frac{\zeta^{[1]}(-k)}{k!}
  &=\frac{\zeta(-k)}{k!}
       \left(\ell-\psi(k+1)
                     -\frac{\zeta^{[1]}(k+1)}{\zeta(k+1)}\right),
                                      &&k\ge1\text{ odd}.
\end{align}
These follow directly by differentiating the Riemann functional equation
at its trivial zeros and at its negative odd arguments.
For $x>1$, use instead the rapidly convergent defining series
\begin{equation}\label{tsc:Li-log-series}
 \Li_0^{[1]}(ze^{-x})=-\sum_{n\ge1}(ze^{-x})^n\log n.
\end{equation}
The checker precomputes the first series coefficients, switches to the
second series after $x=1$, and inserts the analytic first derivative at
the removable endpoint of the subtracted Mellin integrand. The degree and
precision choices are numerical safeguards; no interval enclosure is
claimed. Six root-of-unity endpoint evaluations are also checked against
the independent finite Gamma formula \eqref{tsc:W-boundary}.

\subsection{Explicit singularity subtraction for the cubic example}
Let $f(x)$ be the integrand in \eqref{tsc:explicit-cubic-eq}, and put
\[
 r_1=-\frac{\psi(1/4)}\pi,\qquad r_2=\frac{\psi(1/2)}\pi,
 \qquad g(x)=f(x)-\frac{r_1}{x-1/4}-\frac{r_2}{x-1/2}.
\]
Then
\begin{equation}\label{tsc:cubic-numeric}
 \PV\int_0^1 f(x)\dd x=\int_0^1g(x)\dd x+r_1\log3.
\end{equation}
The removable interior values used in the implementation are
\begin{align*}
 g(1/4)&=-\frac{\psi'(1/4)}\pi-2\psi(1/4)
                         -\frac{r_2}{1/4-1/2},\\
 g(1/2)&=\frac{\psi'(1/2)}\pi-2\psi(1/2)
                         -\frac{r_1}{1/2-1/4},\\
 g(0)&=-\pi+4r_1+2r_2.
\end{align*}
Thus the numerical check uses an ordinary integral with explicitly removed
poles, rather than an extrapolated cutoff sequence.

\subsection{Finite tails and an implementation safeguard}
For the positive-integer Tornheim diagnostics, the Mellin integral is
truncated at a finite $R$. For $r,s\ge0$ integral, unit colors, and
$x\ge\log2$,
\[
 |\Li_r(ze^{-x})\Li_s(we^{-x})|\le4e^{-2x}.
\]
Consequently the omitted tail of \eqref{tsc:T-mellin}, for integer $t\ge1$,
is at most $4\Gamma(t,2R)/(2^t\Gamma(t))$. This finite-tail implementation
also avoids enormous-exponent complex logarithms that can arise when a
black-box infinite-range quadrature samples astronomically large $x$.
It is a safeguard in the newly written numerical checker, not a correction
to an existing repository theorem.

\subsection{Artifact map}
The principal reproducibility files are
\begin{center}
\begin{tabular}{p{0.43\textwidth}p{0.49\textwidth}}
\toprule
File & Purpose\\\midrule
\src{article.tex}, \src{sections/} & Modular complete proof text.\\
\src{article_standalone.tex} & Flattened single-file LaTeX source.\\
\src{article.pdf} & Compiled article.\\
\src{code/verify_exact.py} & 621 exact assertions and coefficient-table generation.\\
\src{code/calculus.py} & Subtracted quadratures and stable polylog jet utilities.\\
\src{code/verify_numerical.py} & 56 numerical comparisons with explicit tolerances.\\
\src{results/} & Executed JSON summaries and finite coefficient tables.\\
\src{notes/} & Source snapshot, audit, proof dependencies, and integration note.\\
\src{verification/} & Build and executed check logs.\\
\src{SHA256SUMS} & Integrity manifest; not a mathematical certificate.\\\bottomrule
\end{tabular}
\end{center}

% END sections/12-computation

% BEGIN references.tex
\begin{thebibliography}{10}
\bibitem{tsc:repo}
V.~Reshetnikov, \emph{ProveIt: Polylogarithms manuscript}, canonical README
and selected integration material, inspected 10 October 2026.
\url{https://github.com/VladimirReshetnikov/ProveIt/tree/main/Analysis/Polylogarithms/docs/manuscript}.
Individual source hashes and read scopes appear in the supplied snapshot.

\bibitem{tsc:shj}
\emph{Shifted Hurwitz jets}, ProveIt research continuation, sections on
canonical finite parts, all-index closure, and further research questions.
Inspected at commit \src{8da6871fc2dd5762467320918ba23edd635bac4f},
blob \src{990369336bdf2794aae3057c3231eed4ecadbef3}.
Repository path:
\src{Analysis/Polylogarithms/docs/reports/stieltjes-correlation-calculus/continuations/shifted-hurwitz-jets/sections.tex}.

\bibitem{tsc:gamma-chapter}
\emph{Log-gamma moments: poles, late coefficients, and effective remainders},
ProveIt canonical manuscript, inspected 10 October 2026.
Repository path:
\src{Analysis/Polylogarithms/docs/manuscript/chapters/07-loggamma-moments.tex}.

\bibitem{tsc:dlmf-zeta}
NIST Digital Library of Mathematical Functions,
\S25.11, \emph{Hurwitz Zeta Function}; especially Eqs.~25.11.3,
25.11.9, 25.11.13, 25.11.17, and 25.11.18.
\url{https://dlmf.nist.gov/25.11}. Accessed 10 October 2026.

\bibitem{tsc:dlmf-poly}
NIST Digital Library of Mathematical Functions,
\S25.12, \emph{Polylogarithms}; especially Eqs.~25.12.10 and 25.12.12.
\url{https://dlmf.nist.gov/25.12}. Accessed 10 October 2026.

\bibitem{tsc:dlmf-gamma}
NIST Digital Library of Mathematical Functions,
\S5.15, \emph{Polygamma Functions}.
\url{https://dlmf.nist.gov/5.15}. Accessed 10 October 2026.

\bibitem{tsc:em1}
O.~R.~Espinosa and V.~H.~Moll,
\emph{On some definite integrals involving the Hurwitz zeta function},
arXiv:math/0012078 (2000).
\url{https://arxiv.org/abs/math/0012078}.

\bibitem{tsc:em2}
O.~R.~Espinosa and V.~H.~Moll,
\emph{On some integrals involving the Hurwitz zeta function: part 2},
arXiv:math/0107082 (2001).
\url{https://arxiv.org/abs/math/0107082}.

\bibitem{tsc:zhao}
J.~Zhao, \emph{A note on colored Tornheim's double series},
arXiv:0907.5106 (2009).
\url{https://arxiv.org/abs/0907.5106}.

\bibitem{tsc:bbb}
D.~H.~Bailey, D.~Borwein, and J.~M.~Borwein,
\emph{On Eulerian log-gamma integrals and Tornheim--Witten zeta functions},
The Ramanujan Journal \textbf{36} (2015), 43--68.
\url{https://doi.org/10.1007/s11139-012-9427-1}.
\end{thebibliography}

% END references

\end{document}
